\documentclass[12pt]{article}

\usepackage[utf8]{inputenc}
\usepackage[T1]{fontenc}
\usepackage{amsmath, amssymb, amsthm}
\usepackage{geometry}
\geometry{a4paper, margin=2.5cm}

\title{Teorema Identității Analitice}
\author{}
\date{}

\theoremstyle{plain}
\newtheorem{theorem}{Teoremă}

\begin{document}

\maketitle

\begin{theorem}
\textbf{Teorema identității analitice.}  
Fie $f$ și $g$ două funcții olomorfe într-un domeniu conex $\Omega \subset \mathbb{C}$.  
Presupunem că există o mulțime $E \subset \Omega$ care are un punct de acumulare în $\Omega$ și pentru care



\[
f(z) = g(z) \quad \text{pentru toate } z \in E.
\]



Atunci



\[
f(z) = g(z) \quad \text{pentru toate } z \in \Omega.
\]



Cu alte cuvinte, două funcții olomorfe care coincid pe o mulțime cu punct de acumulare coincid pe întreg domeniul.
\end{theorem}

\end{document}
